% Regression: p3_eq12 — renders the closed-lasso network defining the L-symbol in equation (12).
% Formula: (L^y_abx)^{z,ij}_{c,k mu} delta_{y,y'} I_y =
% p3 — arXiv:2203.12563 eq (12): the L-symbol definition (closed lasso)
%   (L^y_abx)^{z,ij}_{c,k mu} delta_{y,y'} I_y =
%   [i splits y' -> (a, z); j splits z -> (b, x); mu fuses (a, b) -> c
%    with b reaching mu through an S-curved cross-wire; k fuses (c, x) -> y]
% Free-tier attempt (the staircase has no grid form): \tnput mpo boxes
%   stand in for the fat bars — no bar skin, no black-vs-gray distinction
%   for the module tensor mu — and \tnjoin[route=curve] draws the S-bend.
%   All wires are black: \tnjoin has no wire-hue key, so the paper's red
%   a, b, c wires are lost.  Bar labels inscribe in the boxes instead of
%   sitting under the bars.
\documentclass[varwidth=19cm, border=6pt]{standalone}
\usepackage{amsmath, amssymb}
\usepackage{tenkz}
\begin{document}

\[
  \left(L^{y}_{abx}\right)^{z,ij}_{c,k\mu}\,\delta_{y,y'}\,\mathbb{I}_y
  \;=\;
  \begin{tenkzfree}
    \tnput[mpo]{i}{(0mm,0mm)}{i}
    \tnput[mpo]{j}{(10mm,-4mm)}{j}
    \tnput[mpo]{m}{(28mm,2mm)}{\mu}
    \tnput[mpo]{k}{(36mm,-2mm)}{k}
    % y' in west, a along the top into mu
    \tnjoin[label={y'}]{-6mm,0mm}{i.west}
    \tnjoin[label=a]{i.north east}{m.north west}
    % the step down to j, then the S-curved cross-wire b into mu
    \tnjoin[label=z]{i.south east}{j.north west}
    \tnjoin[route=curve, out=0, in=180, label=b]{j.east}{m.south west}
    % x underneath, c between mu and k, y out east
    \tnjoin[label=x]{j.south east}{k.south west}
    \tnjoin[label=c]{m.east}{k.north west}
    \tnjoin[label=y]{k.east}{42mm,-2mm}
  \end{tenkzfree}
\]

\end{document}
